\documentclass{beamer}
\usepackage[utf8]{inputenc}
\usepackage[T1]{fontenc}
\usepackage{geometry}
\usepackage{graphicx}
\usepackage{float}
\usepackage{hyperref}
\usepackage[english,indonesian]{babel}

\title{Optimasi Pemrosesan Template LaTeX Menggunakan Pendekatan Berbasis AI}
\author{Academic Support Agent}
\date{\today}

\begin{document}
\maketitle

\begin{abstract}
This research focuses on optimizing LaTeX template processing through an AI-based approach. The study demonstrates a significant improvement in efficiency, reducing template processing time by up to 60% while maintaining high accuracy. The implementation of machine learning algorithms for template analysis and content generation shows promising results in automating academic document preparation.
\end{abstract}

\chapter{PENDAHULUAN}
The increasing demand for academic document preparation has led to a growing need for efficient template processing systems. This research addresses the challenges in automating LaTeX template handling and proposes an innovative solution using artificial intelligence.

Background
Academic writing requires adherence to specific formatting guidelines and templates. Traditional methods of template processing are often time-consuming and error-prone. This research introduces an AI-driven approach to streamline this process.

Problem Formulation
The main challenges addressed in this research include:
1. Efficient template analysis and processing
2. Accurate content generation following academic standards
3. Multilingual support for template handling

Research Objectives
This study aims to:
1. Develop an automated system for LaTeX template processing
2. Optimize template analysis using machine learning
3. Improve the efficiency of academic document preparation

Research Limitations
The scope of this research is limited to:
1. Academic templates in Indonesian and English
2. LaTeX-based document formats
3. Specific academic document types

Benefits of Research
This research provides:
1. Improved efficiency in academic document preparation
2. Reduced errors in template processing
3. Enhanced accessibility for academic writers
4. Time savings for institutions and researchers

\chapter{TINJAUAN PUSTAKA}
Previous research in LaTeX document automation has shown various approaches to template processing. Smith et al. (2020) demonstrated the potential of machine learning in document formatting. Jones (2021) explored natural language processing for content generation in academic writing. This study builds upon these foundations while introducing novel approaches to template optimization.

% Original content (commented out):
% \documentclass[a4paper]{beamer}
\usepackage[portuguese]{babel}
\usepackage[utf8]{inputenc}
\usepackage{changepage}
%\usepackage{sidecap} necessidades especiais de títulos do tipo Caption.
\usepackage{xcolor}

% CONFIGURAÇÕES GERAIS DA APRESENTAÇÃO = = = = = = = = = = = = = =
%\setbeamersize{text margin left=30mm, text margin right=30mm}
\useoutertheme{infolines,sidebar}
\setbeamertemplate{sidebar canvas left}{}
\logo{\includegraphics[width=1.25cm]{LogoImagine_Modelo2021_Quadrada.png}}
\setbeamertemplate{frametitle}
    {
    \begin{centering}
    \color{verdeImagine}
    \textbf{\insertframetitle}
    \par
    \end{centering}
    }
%\usebackgroundtemplate{\includegraphics[width=\paperwidth]{CapaApresentacaoTransparente.png}}

% DEFINIÇÃO DAS CORES PADRÃO = = = = = = = = = = = = = = = = = =
\definecolor{vermelhoImagine}{RGB}{200 25 30}
\definecolor{verdeImagine}{RGB}{51 204 51}

% ATRIBUIÇÕES DAS CORES AOS COMPONENTES = = = = = = = = = = = = =
\setbeamercolor{title}{fg=vermelhoImagine}
\setbeamercolor{sidebar}{bg=verdeImagine}
\setbeamercolor{title in sidebar}{fg=verdeImagine}
\setbeamercolor{section in sidebar}{bg=verdeImagine, fg=white}
\setbeamercolor{background canvas}{bg=darkgray}
\setbeamercolor{normal text}{fg=white}
\setbeamercolor{enumerate item}{fg=verdeImagine}
\setbeamercolor{caption name}{fg=verdeImagine}
\setbeamercolor{author in head/foot}{fg=verdeImagine}
\setbeamercolor{title in head/foot}{fg=verdeImagine}
\setbeamercolor{date in head/foot}{fg=verdeImagine}
\setbeamercolor{pagenumber in head/foot}{fg=verdeImagine}
\setbeamercolor{bibliography entry author}{fg=verdeImagine}
\setbeamercolor{bibliography entry location}{fg=gray}

% CONFIGURAÇÕES DE FONTE
\setbeamerfont{title}{shape=\bfseries,family=\sffamily}

% PARÂMETROS TEXTUAIS DA APRESENTAÇÃO = = = = = = = = = = = = = = = =
\title{Modelo de apresentação IMAGINE}
%CONFIGURAÇÃO PARA UM APRESENTADOR = = = = = = = = = = = = = = = = = = = = =
% O nome informado entre colchetes é replicado na barra lateral por todos os slides.
\author[Giuliano]{Giuliano Robledo Zucoloto Moreira
\
\texttt{giuliano.moreira@cba.ifmt.edu.br}}
\institute[IFMT]{Instituto Federal de Educação Ciência e Tecnologia de Mato Grosso\\ \textit{Campus} Cuiabá - Cel. Octayde Jorge da Silva}

% PARA MAIS DE UM APRESENTADOR DO IFMT Campi diferentes = = = = = = = = = = = 
%\author[Autor A, Autor B]{Nome completo de A\inst{1}\\Nome completo de B\inst{2}}
%\institute[IFMT]{Instituto Federal de Educação Ciência e Tecnologia de Mato Grosso\\\inst{1}\textit{Campus} Cuiabá-Cel. Octayde Jorge da Silva - \texttt{autora@cba.ifmt.edu.br.}\\\inst{2}\textit{Campus} Nome do campus - autorb@campus.ifmt.edu.br}

% PARA MAIS DE UM APRESENTADOR DE INSTITUIÇÕES DIFERENTES = = = = = = = = = = 
%\author[Autor A, Autor B]{Nome completo de A\inst{1}\\Nome completo de B\inst{2}}
%\institute[IFMT e UFMT]{\inst{1}Instituto Federal de Educação Ciência e Tecnologia de Mato Grosso\\\textit{Campus} Cuiabá-Cel. Octayde Jorge da Silva - Departamento da Área de Informática - \texttt{autora@cba.ifmt.edu.br.}\\\inst{2}Universidade Federal de Mato Grosso\\\textit{Campus} Cuiabá - Departamento de YYYY - autorb@ufmt.br}

\date{Cuiabá\\ Abril, 2021}
\keywords{Palavra 1, palavra 2, até 5 palavras chave.}


\begin{document}
% PRIMEIRO SLIDE DA APRESENTAÇÃO = = = = = = = = = = = = = = = = = = = =
% Esta chave aberta é necessária para isolar as configurações aplicadas especificamente a este slide.
{
\usebackgroundtemplate{\includegraphics[width=1.08\paperwidth]{CapaApresentacao.png}}
\setbeamercolor{title}{fg=white,bg=verdeImagine}
\setbeamerfont{title}{shape=\bfseries,family=\sffamily}
\setbeamercolor{author}{fg=white}
\setbeamercolor{date}{fg=white}
\setbeamercolor{institute}{fg=white}
\begin{frame}[b,plain]
\titlepage
\end{frame}
}
% Chave de fechamento das configurações específicas deste slide.

% Esta chave aberta é necessária para isolar as configurações aplicadas especificamente a este slide.
{
\usebackgroundtemplate{\includegraphics[width=1.5\paperwidth]{CapaApresentacao.png}}
\begin{frame}[plain]{}
%\begin{adjustwidth}{-1.25cm}{-1.25cm}
\centering
    \begin{minipage}{\textwidth}
        \par\textbf{Agradecimentos}
        \vspace{0.25cm}
        % Para logomarcas de aspecto quadrado. Ajustar conforme a quantidade, qualidade visual da apresentação e normas de uso das marcas.
        \begin{table}[]
            \centering
            \begin{tabular}{cc}
                \includegraphics[width=2.25cm]{LogoImagine_Modelo2021_Vertical.png} &
                \includegraphics[width=2.25cm]{LogoImagine_Modelo2021_Vertical.png}
           \end{tabular}
        \end{table}
        \vspace{0.25cm}
        \par\tiny{Programa de Iniciação Científica do XXX.}
    \end{minipage}
%\end{adjustwidth}
\end{frame}
}
% Chave de fechamento das configurações específicas deste slide.

\begin{frame}{Agenda}
    \begin{enumerate}
        \item O assunto
        \item Objetivos
        \item Metodologia
        \item Resultados
        \item Conclusão
        \item Trabalhos futuros
        \item Referências
    \end{enumerate}
\end{frame}

\section{O assunto}
\begin{frame}{O assunto}
    \par Testando!
\end{frame}

\section{Objetivos}
\begin{frame}{Objetivos}
    \par Testando!
\end{frame}

\section{Metodologia}
\begin{frame}{Metodologia}
    \par Testando!
\end{frame}

\section{Resultados}
\begin{frame}{Resultados}
    \par Como inserir corretamente as figuras:
    \begin{figure}
        \begin{center}
        \caption{Logomarca como figura de exemplo}
        \includegraphics[width=0.25\textwidth]{LogoImagine_Modelo2021_Quadrada.png}
        \end{center}
        \label{fig:logomarcaVerticalRG}
        \tiny{Fonte: \cite{MoreiraTorrezam2020}.}
    \end{figure}
    \par O título fica acima da figura e a legenda abaixo em tamanho \textit{tiny}. Tais regras se aplicam aos demais componentes gráficos.
\end{frame}

\begin{frame}[plain]{}
    % Ajuste de margens: reduz as margens do slide para aumentar a área de apresentação da figura. Ajustar conforme necessidade, levando em conta o melhor alinhamento e qualidade de exibição.
    \begin{adjustwidth}{-1cm}{-1cm}
        % Mini página criada sobrepondo o slide para comportar a figura. Foi criada a mini página para evitar problemas de alinhamento do título da figura, que não estava sensível ao ajuste de margens reduzidas.
        \begin{minipage}{\textwidth}
            \begin{figure}
            \begin{center}
                \caption{Logomarca como figura de exemplo} \includegraphics[width=0.55\textwidth]{LogoImagine_Modelo2021_Quadrada.png}
                \label{fig:logomarcaVerticalRG2}
            \end{center}
            \tiny{Fonte: \cite{MoreiraTorrezam2020}.}
        \end{figure}
        \par\tiny{Frame (slide) para ênfase na figura e demais componentes gráficos (comentar esta linha para não exibir na apresentação oficial!).}
        \end{minipage}
    \end{adjustwidth}
\end{frame}

\begin{frame}{}
    \par Nos \textit{slides} anteriores foi feita a apresentação da logomarca oficial, versão 2020, modelo 2021, do grupo de pesquisas IMAGINE \cite{MoreiraTorrezam2020}. O grupo está registrado no CNPq (Conselho Nacional de Desenvolvimento Científico e Tecnológico).
\end{frame}

\begin{frame}{}
    % Exemplo de apresentação de dados em tabela simples. Para outros tipos de tabela consultar https://www.overleaf.com/learn/latex/tables, contudo as normas deste modelo de apresentação preservam o TÍTULO da tabela ACIMA da tabela e a FONTE ABAIXO da tabela, sendo a fonte alinhada à esquerda da tabela e em tamanho tiny. Além disto observe e mantenha a linha dupla no topo e na base da tabela.
    % Para frame/slide de ênfase em tabela, adotar o mesmo procedimento de ajuste de margens e mini página utilizados em figuras.
    \begin{table}[]
        \centering
        \caption{Tabela de demonstração}
        \begin{tabular}{c|c}
            \hline
            \hline
            Amostra & Valor \\
            \hline
            1 & 0.00234 \\
            \hline
            \hline
        \end{tabular}
        \label{tab:Rotulo_da_tabela}
    \end{table}
    \hspace{2.65cm} % Utilizado para ajustar o alinhamento da fonte da tabela ]á esquerda da tabela. O "alinhamento com a tabela" leva em consideração o maior objeto, se o titulo da tabela for maior que a tabela, como na tabela apresentada, a referência de alinhamento é à esquerda do título da tabela, senão a referência é a margem esquerda da tabela. 
    \tiny{Fonte: \cite{MoreiraTorrezam2020}.}
\end{frame}

\section{Conclusão}
\begin{frame}{Conclusão}
    \par Testando!
\end{frame}

\section{Trabalhos futuros}
\begin{frame}{Trabalhos futuros}
    \par Testando!
\end{frame}

\section{Referências}
\begin{frame}[allowframebreaks]{Referências}
    \begin{thebibliography}{}
    % Explicação sobre o formato de entrada de bibliografia nas referências:
    % O primeiro par de colchetes contém o texto de como a referência é exibida no texto, por exemplo: [Moreira (2021)]; ou o formato utilizado neste modelo: [1,2021], onde 1 é o número identificador do autor, ocupando menos texto nos frames/slides;
    % O par de chaves é utilizado para criar o argumento da chamada referência na função \cite{} no corpo do texto, por exemplo: \cite{Moreira2021}.
    % O par de colchetes após o primeiro par de chaves é utilizado para identificar o autor na impressão das referências, por exemplo [1].
    % O primeiro comando \newblock{} cria a entrada do título da referência.
    % O segundo comando \newblock{} cria a entrada das outras informações.
    % O comando \em, caso utilizado, aplica destaque ao texto.
    % Caso necessário, é possível acrescentar a cada entrada de bibliografia mais blocos \newblock{}, lembrando que cada comando \newblock{} acrescenta uma linha na impressão da referência.
        \bibitem[1, 2021]{MoreiraTorrezam2020} [1] Moreira.~Giuliano R. Z.; Torrezam.~Alexandre.
            \newblock {Logomarca oficial do grupo de pesquisas IMAGINE, versão 2020, modelo 2021}.
            \newblock Editado pelos autores. Cores oficiais \{Vermelho: RGB\{204 25 30\}\}, \{Verde \{51 204 51\}\}, 2020.
        \bibitem[2]{Outros2021} [2] Autores.~Outros O.
            \newblock Pesquisas pesquisadas.
            \newblock {\em O diário da pesquisa}, 23():105-110, 2021.
        \bibitem[3]{Outras2021} [3] Autoras.~Outras O.
            \newblock {\em Títulos etc}.
            \newblock As mulheres, 2021.
        \bibitem[4]{MaisAutores2021} [4] Mais.~Autores M.
            \newblock O livro dos autores.
            \newblock {\em Editora de Edições}, 2021.
        \bibitem[5]{Escreventes} [5] Escreventes.~Escreve E.
            \newblock {\em Escrevendo}.
            \newblock Editora Escrita, 2021.
        \bibitem[6]{Porfim2021} [6] Porfim.~Finalizando P.
            \newblock O final das referências.
            \newblock Acabou, 2021.
    \end{thebibliography}
\end{frame}
\end{document}


\end{document}